\begin{table}[t]
\centering
\caption{Geometric difference $g(K_C\Vert K_Q)$, minimised over classical bandwidth
($N=800$, $\lambda=10^{-2}$, $\sqrt{N}\approx 28.3$). A small $g$ means the
classical kernel can learn anything the quantum kernel can. The last column gives the angle-encoded
fidelity kernel against its best RBF. Rank-deficient kernels are excluded. Generated from the result
files.}
\label{tab:geom}
\footnotesize
\renewcommand{\arraystretch}{1.15}
\setlength{\tabcolsep}{3pt}
\begin{tabular}{@{}llcccc@{}}
\toprule
& & \multicolumn{3}{c}{\textbf{Classical family, vs IQP kernel}} & \textbf{Angle fidelity}\\
\cmidrule(lr){3-5}\cmidrule(lr){6-6}
\textbf{Dataset} & \textbf{IQP kernel} & \textbf{RBF raw} & \textbf{Trig} & \textbf{Phase kernel} & \textbf{vs RBF}\\
\midrule
NSL-KDD & projected & 8.41 & 8.33 & \textbf{2.27} & \multirow{2}{*}{1.07}\\
 & fidelity & 4.91 & 4.88 & \textbf{1.77} & \\
UNSW-NB15 & projected & 12.21 & 12.10 & \textbf{2.44} & \multirow{2}{*}{1.07}\\
 & fidelity & 9.52 & 9.34 & \textbf{1.50} & \\
CICIDS2017 & projected & 8.65 & 8.48 & \textbf{1.95} & \multirow{2}{*}{1.06}\\
 & fidelity & 7.04 & 6.86 & \textbf{1.54} & \\
NF-ToN-IoT-v2 & projected & 9.25 & 9.07 & \textbf{2.23} & \multirow{2}{*}{1.07}\\
 & fidelity & 7.84 & 7.64 & \textbf{1.39} & \\
\bottomrule
\end{tabular}
\end{table}
